\chapter{The Watt: Energy as the Fourth Element}
\label{ch:energy}

The three elements of Part~II share a silent denominator. Every token is ultimately a quantity of electricity passed through silicon, and the AI war's longest-lead-time input is neither the model, the chip, nor the fab---it is the watt.

This chapter states the asymmetry carefully, because the obvious comparison is the wrong one. \emph{Nameplate capacity added is not deliverable power at a computing site}, and a book that lets a solar gigawatt stand in for a gigawatt of 24/7 data-centre supply has made exactly the category error it accuses others of making elsewhere. Table~\ref{tab:firmpower} therefore replaces the single headline figure with a four-row ledger, and the argument that follows rests on the rows that survive it: not that China adds ten times more nominal capacity, but that it combines much faster buildout with long-distance transmission, state-directed siting, retained dispatchable firming, and---decisively---the ability to attach large new loads on a schedule the American interconnection queue cannot match.

\section{The supply asymmetry}

The 2025 numbers, stated at their correct denominators [V]. China's national energy administration reports \emph{452\,GW of new renewable capacity} in 2025---318\,GW solar and 120\,GW wind---within total additions of roughly 540\,GW, bringing installed capacity to 3.89\,TW at year-end and about 4.04\,TW by mid-2026, on $\sim$\$500B of annual energy investment. The US Energy Information Administration reports 53\,GW of new utility-scale capacity in 2025, including storage---a record American year~\cite{china-capacity,us-additions,nea-2025}.

\begin{table}[htbp]
  \centering
  \small
  \caption{Four denominators, not one. A gigawatt of solar nameplate is not a gigawatt of firm data-centre supply; the strategic claim must be built from the rows that survive the distinction.}
  \label{tab:firmpower}
  \begin{tabular}{p{40mm}p{50mm}p{52mm}}
    \toprule
    Metric & China & United States \\
    \midrule
    New nameplate capacity (2025) & $\sim$540\,GW total; 452\,GW renewable (318 solar, 120 wind) [V] & 53\,GW utility-scale incl.\ storage [V] \\
    \addlinespace
    Expected annual generation & capacity-factor adjusted; generation 10,573\,TWh, $+$503\,TWh y/y [V] & 4,430\,TWh, growth modest [V] \\
    \addlinespace
    Firm / dispatchable contribution & record 95\,GW coal commissioned as flex; $\sim$66\,GW nuclear toward 110\,GW by 2030; two-thirds of world grid storage [V/R] & gas turbines sold out to $\sim$2030; $<$1.7\,GW of marquee nuclear restarts [V] \\
    \addlinespace
    Deliverable data-centre headroom & hub-specific: 45 UHV lines, four new 8-GW UHVDC corridors in 2025; state-directed siting into curtailed-renewable regions [V/P] & balancing-area specific: $\sim$2,060\,GW queue, $>$3-yr median to agreement; 12+ states with moratorium bills [V] \\
    \bottomrule
  \end{tabular}
\end{table}

Three qualifications belong with the table. First, storage is not primary generation, and a national aggregate surplus does not imply firm deliverability at any particular computing hub---the correct unit of analysis is the balancing area or the provincial grid, not the country. Second, falling average generator utilization hours are consistent with several stories besides surplus: rapid renewable additions, curtailment, congestion, weak demand, or displacement of thermal generation. The claim that the marginal Chinese AI data centre bids into genuine headroom therefore requires \emph{regional} evidence---average and marginal industrial tariffs, curtailment and congestion, firm generation and storage, transmission import capacity, and time-to-interconnect a 100--500\,MW load---for Inner Mongolia, Gansu, Guizhou, Guangdong, and the Yangtze River Delta specifically. That evidence is assembled unevenly in public sources and is the single largest gap in this chapter [A]. Third, global context: the IEA's projection of $\sim$945\,TWh of data-centre consumption in 2030 remains under 3\% of world electricity---\emph{the stress is regional and temporal, not globally energy-limited}~\cite{iea-energy-ai}. Nothing in this chapter should be read as a claim that the planet runs out of electricity; the claim is that particular buildouts, in particular interconnection queues, on particular schedules, do. China's generation ran 10,573\,TWh in 2025 against America's 4,430---2.4$\times$---and its \emph{annual demand growth} ($+$503\,TWh) now exceeds Germany's total consumption~\cite{ember-2026}. The composition compounds the point (wind plus solar now 1.84\,TW, nearly half the fleet, with clean generation passing 40\% of output in H1 2026); ten-reactor-per-year nuclear approvals toward 110\,GW by 2030; a record 95\,GW of coal commissioned as dispatchable flex; and roughly two-thirds of the world's grid-battery installations---one Chinese \emph{month} of BESS deployment in 2025 exceeded the US full year~\cite{china-capacity,nuclear-china,coal-2025,bess-2025}. Around it, the delivery machine: dozens of UHV lines moving western generation eastward, four new 8-GW UHVDC corridors entering service in 2025 alone, and roughly 28 more lines planned by 2030~\cite{uhv,oies-china-dc}. Capacity is running ahead of demand---average generator utilization hours are \emph{falling}~\cite{china-capacity}---but the natural conclusion from this overreaches, and the disciplined statement is the one that belongs on the record: \emph{the national data are consistent with aggregate surplus, but do not establish firm, deliverable surplus at a specific computing hub}. The distinction is not pedantic: 2025 national wind/solar curtailment ran 5.4\%, breaching the official 5\% target, and Q1 2026 saw large weather-adjusted capacity-factor declines concentrated in exactly the hub provinces (Inner Mongolia foremost)---driven less by missing wires than by inflexible coal contracts and annual provincial trading, which is a \emph{market-design} constraint on deliverability that a data center experiences just as physically as a congestion constraint~\cite{doc136}. The busbar casebook of Section~\ref{sec:busbar-casebook} is where the hub-level evidence now lives.

\begin{figure}[htbp]
\centering
\begin{tikzpicture}
\begin{axis}[bookbar, height=5.0cm, bar width=13pt,
  ylabel={Generating capacity added (GW)},
  title={Annual additions to generating capacity},
  symbolic x coords={2023,2024,2025},
  xtick=data, ymin=0, ymax=620,
  nodes near coords, nodes near coords style={font=\scriptsize\color{inkSec}},
  legend style={at={(0.02,0.98)}, anchor=north west},
]
\addplot[fill=sOne] coordinates {(2023,300) (2024,429) (2025,540)};
\addlegendentry{China}
\addplot[fill=sTwo] coordinates {(2023,32) (2024,45) (2025,53)};
\addlegendentry{United States}
\end{axis}
\end{tikzpicture}

\vspace{1mm}

\begin{tikzpicture}
\begin{axis}[bookbar, height=4.2cm, width=0.72\linewidth, bar width=16pt,
  ylabel={GW},
  title={The queue and the ask, for scale},
  symbolic x coords={{US 2025 built},{OpenAI's annual ask},{US interconnection queue}},
  xtick=data, ymin=0, ymax=2400,
  xticklabel style={font=\scriptsize\color{inkSec}, align=center},
  nodes near coords, nodes near coords style={font=\scriptsize\color{inkSec}},
]
\addplot[fill=sFour] coordinates
  {({US 2025 built},53) ({OpenAI's annual ask},100) ({US interconnection queue},2060)};
\end{axis}
\end{tikzpicture}
\caption[The energy asymmetry]{The fourth element. \emph{Above}: annual additions to generating capacity, China
against the United States---China added roughly ten times as much in 2025, into falling utilization hours, while
the US record year was 53\,GW [V]~\cite{china-capacity,us-additions}. \emph{Below}: the American constraint stated
three ways---what was built in the best year on record, the 100\,GW/year one AI laboratory told the White House
the country needs, and the $\sim$2,060\,GW sitting in the interconnection queue at median waits over three years
[V/P]~\cite{us-queue,openai-100gw}. One side's binding constraint is capital; the other's is physics and
permitting, and permitting is the harder one to finance away.}
\label{fig:energy}
\end{figure}


The United States is the mirror image: demand arriving into scarcity. Here a category error runs through much of the public debate, and disentangling it sharpens the argument rather than weakening it. The famous $\sim$2,000-GW-plus interconnection queue is a queue of \emph{proposed generation and storage projects} seeking transmission interconnection---roughly 2,290\,GW active at the end of 2024, with median request-to-operation times above four years and only 13\% of 2000--2019 requested capacity ever built~\cite{lbl-queued}. It says nothing directly about how long a 100--500\,MW \emph{data-center load} waits. The US constraint properly decomposes into five queues that compound: (1) the generation-side queue above, which throttles the supply additions needed to serve new load; (2) \emph{large-load connection} studies and agreements, for which no comprehensive national statistic exists---the definitive study of the problem states plainly that no cross-utility survey has been done---but where the observable instances are grim: utility load studies of 6--18 months followed by 18--60 months of construction where line extensions are needed, and Northern Virginia transmission-level hookups quoted at up to seven years against 18--24-month shell construction~\cite{lbl-speed,dominion-wait}; (3) transmission expansion, on 7--10-year planning-to-energization cycles; (4) equipment---gas turbines effectively sold out into 2029--2030, large power transformers at 24-to-60-month lead times; and (5) resource adequacy, now priced: PJM's capacity auctions have cleared at or near their caps three years running, \$16.4B for 2027/28 and again for 2028/29, with the RTO missing its own reliability requirement in both~\cite{gas-turbines,transformers,pjm,pjm-2829}. Data-center load, 62\,GW in 2025, is projected to more than double by 2030, driving 90 of the 166\,GW of forecast five-year peak-load growth~\cite{sp-dc,grid-strategies}; PJM's capacity auctions have set three consecutive records (\$16.4B for 2027/28, $\sim$94\% of load growth attributed to data centers), ratepayer politics has produced data-center-moratorium bills in at least twelve states, and the physical inputs are queued years deep---gas turbines effectively sold out to 2030 (GE Vernova's backlog at 116\,GW, taking reservations for 2031), large transformers at two-year-plus lead times and prices up 70\% since 2019~\cite{pjm,moratoria,gas-turbines,transformers}. The marquee nuclear responses---restarting Three Mile Island for Microsoft (2027 target) and Palisades (slipping)---total under 1.7\,GW: measurement error against the Chinese fleet~\cite{tmi-palisades}. The system's own principals have said the quiet part: Microsoft's CEO---``it's not a supply issue of chips; it's actually the fact that I don't have warm shells to plug into''---and OpenAI's letter to the White House demanding \emph{100\,GW per year} of new US capacity, justified explicitly by China's 429-GW year and an ``electron gap''~\cite{nadella-power,openai-100gw}. The ask is nearly double the best year the US grid has ever had.

\section{Energy as industrial policy: the token factory's power bill}

The asymmetry is not merely quantitative; it is being \emph{aimed}. The East-Data-West-Computing program placed eight national computing hubs where the curtailed renewables are; 2024 policy requires 80\% renewable sourcing for new hub data centers; and in November 2025---in the same week Beijing barred foreign AI chips from state-funded data centers---provinces including Gansu, Guizhou, and Inner Mongolia cut electricity prices \emph{up to 50\%, to $\sim$RMB\,0.4 (\$0.056)/kWh, exclusively for facilities running approved domestic chips}~\cite{edwc,china-dc-power,datacenter-mandate}. That single instrument welds three elements of this book together: it converts the grid surplus into a subsidy, the subsidy into demand for Ascend/Cambricon silicon, and the silicon into tokens priced below Western marginal cost. It also directly attacks the domestic chips' principal weakness---30--50\% worse energy efficiency than NVIDIA equivalents---though the correct description is not that the wasted watts become free. Extra watts still require generation, transmission, transformers, cooling, water, and backup capacity, and they still carry an opportunity cost; what the tariff does is \emph{socialize part of the operating-cost penalty and accelerate domestic-chip utilization}, transferring the efficiency gap from the buyer's income statement to the provincial grid's~\cite{china-dc-power}. Jensen Huang's London formulation---that for his Chinese competitors ``power is free''---is a useful exaggeration of a real subsidy, not a description of the physics~\cite{huang-quotes}.

\subsection{The value of flexible computing}
\label{sec:flex}

One further asymmetry is under-analysed and probably favours the state-directed system. AI workloads are not uniformly time-bound: interactive serving must run now, but training, batch inference, synthetic-data generation, and much of the scientific execution layer of Chapter~\ref{ch:sizing} can be interrupted, checkpointed, migrated geographically, or scheduled to follow renewable output. The economic value of that flexibility---interruptible training, checkpoint-aware demand response, geographic workload migration, renewable-following batch, and the substitution of \emph{compute} storage (precomputed tokens, caches, distilled models) for \emph{energy} storage---accrues to whoever can coordinate siting, transmission, and scheduling across a national footprint. A network of eight state-designated computing hubs wired to curtailed western generation by dedicated UHV corridors is, in effect, a demand-response instrument at national scale; a set of independently sited private campuses bidding into congested regional markets is not. If this chapter has a claim that deserves more attention than the headline gigawatts, it is this one [S].

Honesty about magnitudes, following the OIES analysis [A]~\cite{oies}: electricity is $\sim$8\% of a US AI data center's three-year TCO---hardware dominates---so cheap power alone does not decide token economics; average Chinese spot prices ($\sim$\$43--58/MWh) overlap US regional averages; and coastal Chinese rates are unremarkable. The Chinese energy advantage is therefore not the tariff but the \emph{headroom}: the ability to attach tens of gigawatts of new load per year without queues, auctions, moratoria, or political backlash, plus the targeting precision of a state that can price discriminate by chip vendor. The US constraint is not generation economics but delivery time: in the scenario logic of Chapter~\ref{ch:trade}, the binding question is which grid can energize capacity on the schedule the capital commitments assume. On current evidence one side's constraint is money and the other's is physics and permitting---and physics-and-permitting is the harder constraint to finance away.

\section{The busbar casebook: ten regions, one table}
\label{sec:busbar-casebook}

The regional evidence named above as this chapter's largest gap ought to be assembled rather than promised, and this section assembles as much of it as the public record supports. Tables~\ref{tab:busbar-us} and~\ref{tab:busbar-cn} apply one standardized frame to five US regions and five Chinese hubs. Grades are per-cell: interconnection statistics and auction prices are [V]; Chinese hub tariffs are [P] (state-media and official reporting); curtailment attributions are [A]. The two systematic gaps, disclosed: no comprehensive US survey of large-load connection times exists (the definitive study says so~\cite{lbl-speed}), and Chinese contracted data-center tariffs below the standard industrial rate are reported for the hubs but not independently auditable.

\begin{table}[htbp]
  \centering
  \scriptsize
  \caption{US regions: what a 100--500\,MW load actually faces (as of this edition's date).}
  \label{tab:busbar-us}
  \setlength{\tabcolsep}{3pt}
  \begin{tabular}{p{20mm}p{30mm}p{28mm}p{32mm}p{28mm}}
    \toprule
    Region & Time to connect & Price signal & Renewables/curtailment & Political durability \\
    \midrule
    N.\ Virginia / PJM (Dominion) & up to $\sim$7\,yr transmission-level hookups; $>$36-mo substation queue vs 18--24-mo shell build [V]~\cite{dominion-wait} & capacity cleared at/near cap 3 yrs running; \$16.4B for 2027/28 and again 2028/29, both short of reliability requirement [V]~\cite{pjm,pjm-2829} & modest curtailment; 2.6\,GW offshore wind due 2026 & Loudoun ended by-right development; moratorium votes pending; new $>$25\,MW rate class with 85\% minimum demand charge from 2027 [V] \\
    \addlinespace
    ERCOT / Texas & fastest connect-and-manage reputation, but large-load queue $>$230\,GW (4$\times$ in a year), only $\sim$7.5\,GW connected; behind-the-meter gas is the workaround [V]~\cite{ercot-queue} & energy-only, scarcity to \$5,000/MWh; EIA projects data-center demand could lift wholesale $\sim$79\% by 2027 & $>$8\,TWh wind+solar curtailed 2024, transmission-bound West & SB6: study fees, deposits, duplicate-request disclosure, and a mandatory remote-curtailment ``kill switch'' for new $\ge$75\,MW loads [V] \\
    \addlinespace
    MISO & load connection via member utilities; tier-2 metros 12--24 mo for moderate loads; ERAS fast-track for adequacy-critical generation & capacity whiplash: \$666.50/MW-day (2025/26) $\to$ \$424.30 (2026/27 summer) & wind-heavy north; curtailment not separately reported & cost-allocation fights (\$280M auction adjustment contested at FERC) \\
    \addlinespace
    SPP & HILL framework (2025): 90-day study track for large loads paired with new generation; provisional and non-firm load service approved & steepest relative growth of any RTO: peak 56\,$\to$\,105\,GW in a decade & wind curtailment up $\sim$6$\times$ 2020--24 & new framework, untested at scale \\
    \addlinespace
    CAISO / PNW & utility-owned process; PG\&E standardized transmission-level tariff (2025); CAISO exploring non-firm ``speed-to-power'' service & California $\sim$26\textcent/kWh all-sector---the cost ceiling of the US map; PNW hydro among the cheapest & CAISO 3.4\,TWh curtailed 2024 (93\% solar); PNW: 9\,GW capacity shortfall by 2030, cut to 0--3\,GW \emph{if} data-center curtailment is allowed & PNW politics of hydro preference; CA affordability backlash \\
    \bottomrule
  \end{tabular}
\end{table}

\begin{table}[htbp]
  \centering
  \scriptsize
  \caption{Chinese hubs: the same frame. Tariffs are contracted data-center rates where reported [P]; standard industrial single-part rates run 0.72--0.81 RMB/kWh across these provinces for comparison.}
  \label{tab:busbar-cn}
  \setlength{\tabcolsep}{3pt}
  \begin{tabular}{p{20mm}p{30mm}p{28mm}p{32mm}p{28mm}}
    \toprule
    Hub & Time to connect & Price signal & Renewables/curtailment & Political durability \\
    \midrule
    Inner Mongolia (Ulanqab/ Hohhot) & hub-designated; grid built ahead of load; gigawatt-class campuses proceeding & 0.32--0.38 RMB/kWh (\$0.045--0.053) data-center band [P]~\cite{hub-tariffs} & $>$80\% green supply claimed; 200+ free-cooling days; \emph{but} among the worst hidden-curtailment provinces (CHP inflexibility) [A]~\cite{doc136} & tariff is policy, not market; Doc-136 early adopter \\
    \addlinespace
    Gansu (Qingyang) & fastest-growing hub: 168k PFLOPS by mid-2026; dedicated 2\,GW green plant built for the hub & 0.398 RMB/kWh (\$0.056), stated as stabilized; 35--45\% below eastern cities [P]~\cite{hub-tariffs} & $>$30\% variable renewables, chronic export dependence; 8\,GW UHVDC to Shandong (2025) & subsidized rate outliving the surplus is the named risk (\S\ref{sec:busbar}) \\
    \addlinespace
    Guizhou (Gui'an) & dedicated substations added within quarters; DC electricity use $+$517\% y/y Q1 2025 [P]~\cite{hub-tariffs} & subsidized interior pricing (historically $\sim$0.35 RMB/kWh deals) [A]~\cite{oies-china-dc} & 35\% renewables + large hydro; mild climate & 30\% AI-adaptation vouchers; hyperscaler-anchored (Huawei, Tencent, Apple) \\
    \addlinespace
    Guangdong (Shaoguan) & GBA hub node; 500k-rack target & coastal 0.60--0.75 RMB/kWh (\$0.08--0.11)---above the Virginia average [A]~\cite{oies-china-dc} & $\sim$70\% coal-sourced eastern power; latency keeps load anyway & Doc-136 finalized late \\
    \addlinespace
    Yangtze Delta / Shanghai & incumbent digital-economy center; East-Data-West-Computing routes batch load out & 0.74--0.80 RMB/kWh industrial; the price the western hubs are arbitraging & $\sim$70\% coal; PUE $\le$1.25 mandate binds hardest here & Jiangsu among last provinces without Doc-136 rules---renewable-pricing uncertainty at the demand center \\
    \bottomrule
  \end{tabular}
\end{table}

Three findings follow from the two tables. \emph{First, the two systems fail at opposite ends of the same pipe.} Every US row's binding constraint is connection time or capacity price---the queue, the auction, the transformer; every Chinese hub row's binding constraint is durability---whether an administered tariff, built on curtailed renewables and a policy decision, outlives the surplus and the political attention that created it. \emph{Second, the intra-China spread is the real story}: 0.32 RMB/kWh in Hohhot against 0.80 in Shanghai is a 2.5$\times$ spread that East-Data-West-Computing exists to arbitrage, and the domestic-chip discount stacks on top of it---the busbar map \emph{is} the industrial policy. \emph{Third, flexibility is worth more than generation on both maps}: the PNW's 9-GW shortfall collapses to 0--3\,GW if data-center load is curtailable, ERCOT now requires a kill switch by law, and CAISO is designing non-firm service---while China's hubs presume interruptible batch workloads by design. The convergence, from opposite directions, on \emph{curtailable compute} as the grid's preferred customer is quiet confirmation of Section~\ref{sec:flex}'s claim that workload flexibility is the under-priced strategic asset of the whole energy element.

\subsection{Busbar-level deliverability: the unit the argument must be priced in}
\label{sec:busbar}

The analysis has to descend one more level, because only at that level do this chapter's national aggregates become a procurement-grade quantity. The decision-relevant unit is neither national generation nor even the balancing area but \emph{firm deliverable megawatts at a specific busbar on a specific date at a specific price}. For any candidate computing site, the ledger is five numbers: (1) firm import plus local dispatchable capacity at the busbar, net of existing load, under the binding seasonal constraint (not the annual average); (2) time-to-interconnect for a 100--500\,MW block, including studies, upgrades, and transformer lead times; (3) the busbar price---tariff plus congestion plus losses plus firming---for a 24/7 load-following profile, which can differ from the zonal average by a factor of two in both directions; (4) curtailment coincidence, i.e.\ whether the cheap hours align with interruptible workload (Section above); and (5) the political durability of all four, which moratorium politics in twelve US states and provincial tariff discrimination in China both demonstrate is a live variable, in opposite directions. The two systems' asymmetry restated at this level: the American constraint binds at (2)---the queue---almost regardless of (1) and (3); the Chinese hub system was \emph{sited} to make (1) and (4) favorable and uses state pricing to set (3), so its binding constraint is (5)'s inverse, the risk that subsidized tariffs outlive the surplus that justified them. The busbar price and time-to-interconnect are precisely the ``energy'' row of the rent-migration table (Table~\ref{tab:rent-migration}); a reader who wants this chapter's thesis as a tradeable indicator should watch those two series, hub by hub, and this book commits to carrying them in the living appendix as regional data becomes assembled.

\section{Where the watt meets the architecture}

The energy element feeds back into the technical argument of Chapter~\ref{ch:compute} in three ways. First, cheap-and-abundant power changes the efficiency frontier: a CloudMatrix-class system at 3.9$\times$ the power of its NVIDIA counterpart is a bad trade in Virginia and an acceptable one in Guizhou---brute-force parity is a rational strategy precisely where watts are the cheap input~\cite{cloudmatrix}. Second, the capacity-first LPDDR architecture's MFU penalty (67\% vs 89\%) is, among other things, an energy trade a surplus grid can afford [M]. Third, the decentralization thesis inverts in energy terms---and the point has to be stated consistently with this book's own power ledger (Table~\ref{tab:128power}), on pain of contradiction: a 128-socket (64-node) sovereign pretraining cluster is modeled at roughly \emph{half a megawatt}, and even larger replicated or multi-rack installations remain in the low-megawatt range---an ordinary industrial feeder, deployable in jurisdictions that could never host a gigawatt campus, which is most of the world. The Global South adoption channel of Chapter~\ref{ch:model} is, in energy terms, a market for machines sized to real grids.

The IEA's trajectory---global data-center consumption from 415\,TWh (2024) toward $\sim$945\,TWh (2030), with the US and China as the two poles~\cite{iea-energy-ai}---frames the endgame: the AI war's operating cost curve is converging on the electricity cost curve. Part~I documented who won the technologies that produce electricity. Solar modules, batteries, UHV transmission, and the coming terawatt of storage are the same Chinese industrial commons that Chapters~\ref{ch:solar}--\ref{ch:ev} described; the watt is where the precedent wars and the AI war become one war. It is the fourth element because it was the first three's prize.
